% <<<GENERATED by generate_spec_twins.py
% Twin of HOUSE_STANDARDS.toml -- machine-generated region below is
% regenerate-only (rerun generate_spec_twins.py --force). Hand-curated
% maths belongs in the marked sections.
\documentclass[11pt]{article}
\usepackage{amsmath,amssymb,mathtools}
\usepackage{booktabs}
\usepackage[margin=1in]{geometry}
\usepackage{hyperref}

% PACKAGE JUSTIFICATION (10_lang_latex requires justification in the header):
%   amsthm         — the guide requires each invariant as a NUMBERED THEOREM;
%                    amsthm supplies the theorem/proof environments for that.
%   newunicodechar — the TOML SSOT is UTF-8 and states its invariants in logic
%                    notation. The GENERATED region reproduces that text verbatim,
%                    so pdflatex must be taught each glyph rather than the SSOT
%                    being de-mathematised to suit the typesetter.
%   fontenc/lmodern — amsthm typesets a proposition body in italic; a \textsc
%                    inside it needs the OT1-undefined small-caps-italic shape.
%                    T1 supplies it and lmodern keeps T1 scalable, not bitmap.
\usepackage[T1]{fontenc}
\usepackage{lmodern}
\usepackage{amsthm}
\usepackage{newunicodechar}
\newtheorem{prop}{Proposition}[section]
\providecommand{\Box}{\square}
\providecommand{\Diamond}{\lozenge}
\newunicodechar{Δ}{\ensuremath{\Delta}}
\newunicodechar{Σ}{\ensuremath{\Sigma}}
\newunicodechar{Φ}{\ensuremath{\Phi}}
\newunicodechar{α}{\ensuremath{\alpha}}
\newunicodechar{δ}{\ensuremath{\delta}}
\newunicodechar{θ}{\ensuremath{\theta}}
\newunicodechar{μ}{\ensuremath{\mu}}
\newunicodechar{ρ}{\ensuremath{\rho}}
\newunicodechar{⁵}{\textsuperscript{5}}
\newunicodechar{⁻}{\textsuperscript{\ensuremath{-}}}
\newunicodechar{₀}{\textsubscript{0}}
\newunicodechar{₁}{\textsubscript{1}}
\newunicodechar{₂}{\textsubscript{2}}
\newunicodechar{↦}{\ensuremath{\mapsto}}
\newunicodechar{⇒}{\ensuremath{\Rightarrow}}
\newunicodechar{∀}{\ensuremath{\forall}}
\newunicodechar{∃}{\ensuremath{\exists}}
\newunicodechar{∅}{\ensuremath{\emptyset}}
\newunicodechar{∈}{\ensuremath{\in}}
\newunicodechar{−}{\ensuremath{-}}
\newunicodechar{∘}{\ensuremath{\circ}}
\newunicodechar{∥}{\ensuremath{\parallel}}
\newunicodechar{∧}{\ensuremath{\wedge}}
\newunicodechar{∩}{\ensuremath{\cap}}
\newunicodechar{≈}{\ensuremath{\approx}}
\newunicodechar{≠}{\ensuremath{\neq}}
\newunicodechar{≡}{\ensuremath{\equiv}}
\newunicodechar{≤}{\ensuremath{\le}}
\newunicodechar{≥}{\ensuremath{\ge}}
\newunicodechar{≫}{\ensuremath{\gg}}
\newunicodechar{⊆}{\ensuremath{\subseteq}}
\newunicodechar{⊇}{\ensuremath{\supseteq}}
\newunicodechar{⊑}{\ensuremath{\sqsubseteq}}
\newunicodechar{⊥}{\ensuremath{\bot}}
\newunicodechar{⊨}{\ensuremath{\models}}
\newunicodechar{⊭}{\ensuremath{\nvDash}}
\newunicodechar{⋀}{\ensuremath{\bigwedge}}
\newunicodechar{⋁}{\ensuremath{\bigvee}}
\newunicodechar{⋃}{\ensuremath{\bigcup}}
\newunicodechar{⌈}{\ensuremath{\lceil}}
\newunicodechar{⌉}{\ensuremath{\rceil}}
\newunicodechar{□}{\ensuremath{\Box}}
\newunicodechar{▸}{\ensuremath{\blacktriangleright}}
\newunicodechar{◇}{\ensuremath{\Diamond}}
\newunicodechar{⟺}{\ensuremath{\Longleftrightarrow}}

\title{HOUSE\_STANDARDS --- Formal Specification}
\author{Beagle paired-artifact doctrine (10\_lang\_latex)}
\date{\today\\ \texttt{toml v15.1\quad sha256:8a81da4cdb515c18727d4ce5be1a29a42e6f2de0d775c7fecdad6ccbec843d9e}}

\begin{document}
\maketitle

\section{Derived Prompt Frame}
The injected frame is derived at runtime from the TOML SSOT:
\begin{equation}
  X = \Lambda(T_{HOUSE\_STANDARDS}, \mathrm{ctx}), \qquad \mathrm{bytes}(X) \le B,
  \label{lbl:derive-frame}
\end{equation}

% <<<GENERATED from TOML -- regenerate only >>>
\section{Rule Coverage (generated)}
\begin{itemize}
  \item \texttt{HS-GATE-01} --- Do not let an unevaluable or skipped check fall through to ALLOW, and do not allowlist an intent name and then trust its parameters. Validate before execute, deny on uncertainty, and emit a ValidationResult with an audit\_id on EVERY path — accepts included. %% lbl:HS-GATE-01
  \item \texttt{HS-GATE-02} --- Do not retry in the transport and do not report loss as one `dropped\_total`. Express admission as a conjunction over locally observable predicates, derive the drop taxonomy by negating it in a declared order, give every reason its own counter, and leave retry to the layer that knows the intent. %% lbl:HS-GATE-02
  \item \texttt{HS-GATE-03} --- Do not give a channel one global backpressure behaviour. Declare a QoS class per topic at creation (LOSSLESS blocks, LOSSY(n) drops and counts, DURABLE tees), default to the class that loses nothing, make eviction direction a separate flag, and parameterize the tests over the truth table. %% lbl:HS-GATE-03
  \item \texttt{HS-GATE-04} --- Do not gate a noisy signal on one threshold, and do not paper over the flapping with debouncing or input smoothing. Declare an entry threshold and a strictly lower exit threshold, hold state inside the band, set band width from measured noise amplitude, and emit the level rather than an event per sample. %% lbl:HS-GATE-04
  \item \texttt{HS-GATE-05} --- Do not treat a config boolean as consent, do not scan on the way out, and do not tell a requester why a consent check failed. Scan before persistence, gate export on signed-consent ∧ clean-scan ∧ namespace-policy failing closed on every row including unevaluable ones, count every reject reason, record every export so revocation tombstones have somewhere to go, and log hashes never payloads. %% lbl:HS-GATE-05
  \item \texttt{HS-GATE-06} --- Do not hand an agent an exit code and expect it to fix the cause. Emit one structured verdict per gate run — PASS/FAIL/SKIPPED, a per-engine table, and per-violation file:line, rule\_id, measured-vs-expected, and owning domain — and cap re-runs per finding set so a refused fix escalates instead of looping. %% lbl:HS-GATE-06
  \item \texttt{HS-GATE-07} --- Do not re-handshake a backend on every call, and do not trust a row you verified once and never rechecked. Verify each sibling once, persist a trust row carrying version AND fingerprint, and re-check the fingerprint before every proxied call — a changed binary is re-verified or refused, never called blind. %% lbl:HS-GATE-07
  \item \texttt{HS-GATE-08} --- Do not install an absent platform's toolchain to make a build configure, and do not hand-write a per-test stub for it. Put the real backend behind a build option defaulting OFF, ship a substitute driver implementing the identical interface from the same repository, run ONE conformance suite against both, and print a capability matrix so substitute-mode is never mistaken for hardware. %% lbl:HS-GATE-08
  \item \texttt{HS-MUT-01} --- Do not compare path strings to prove containment, do not follow a symlink out of a declared root, and do not let discovery delete. Resolve to real paths and test the descendant relation, refuse escapes and report the refusal as its own outcome, separate scan/adopt/rotate, simulate by default so deletion needs an explicit force, bound each step with a timeout, and discard an inherited config's unsafe defaults. %% lbl:HS-MUT-01
  \item \texttt{HS-MUT-02} --- Do not perform an effect and then log it, and do not compute the inverse at revert time. Commit an append-only hash-chained entry carrying intent AND the inverse — computed while the pre-state is still observable — before touching the world; make inverses idempotent and leaf-classified; alert on a chain break; and test by killing the process at every step boundary. %% lbl:HS-MUT-02
  \item \texttt{HS-MUT-03} --- Do not mutate a read-mostly index in place, do not lock readers for an update that almost never happens, and do not re-read the snapshot pointer per batch element. Build a new immutable snapshot off the hot path, publish it with one atomic pointer store, pin ONE snapshot for a whole operation, free the old one by reference count rather than at swap time, give each snapshot an id, and soak-test for flat memory beside the correctness drill. %% lbl:HS-MUT-03
  \item \texttt{HS-MUT-04} --- Do not make the ring infer a message length, do not hold the lock across the caller's copy, and do not crash the writer on a corrupt ring. Reserve a slot, hand back a writable view, take the length from the caller on commit, advance the head only on commit, default a forgotten commit to an idempotent zero, health-check before every operation, and quarantine a corrupted ring for operator re-provisioning rather than retrying. %% lbl:HS-MUT-04
  \item \texttt{HS-MUT-05} --- Do not hand a raw pointer into a recyclable ring slot, and do not let a SIMD parser read to the exact end of a payload. Mint a view guard capturing a monotonic slot epoch, validate the epoch on every dereference and raise on mismatch, keep consumer views read-only, and pad every slot with a full vector width of trailing bytes. %% lbl:HS-MUT-05
  \item \texttt{HS-BOUND-01} --- Do not degrade a missing optional backend to a no-op stub or to a bare None. Declare the capability as a Protocol, return Present(backend) | Absent(capability, reason) as the public union so the type checker refuses every undiscriminated call site, keep Absent falsy but distinct from None, leave the fallback policy with the caller who owns it, and print the resolved capability set in the doctor surface. %% lbl:HS-BOUND-01
  \item \texttt{HS-BOUND-02} --- Do not move code and tidy it in the same commit, and do not leave a shim that re-implements. Record a characterization suite against the module in place FIRST, move it unchanged, leave a facade that only re-exports, run the suite through both paths for the whole deprecation window, freeze the public surface as a version-controlled file asserted by a snapshot test, state the removal release in the shim itself, and extract in dependency order. %% lbl:HS-BOUND-02
  \item \texttt{HS-BOUND-03} --- Do not require a host-core edit to register an extension. Publish plugins as packages under a PEP 420 namespace at <namespace>.plugins.<name>, have the host scan the namespace and PLUGINS\_DIR and import each with a per-plugin timeout, require `app` and treat the lifecycle hooks (register\_services, register\_routes, get\_env\_vars, shutdown) as optional, ship config templates as package-data, and keep cli.py thin. %% lbl:HS-BOUND-03
  \item \texttt{HS-BOUND-04} --- Do not build a CLI, an MCP server, and a TUI as three implementations. Put run/format/config in one core and make each surface a thin adapter, sharing one config schema, one totally ordered precedence rule (CLI args > project TOML > user config > defaults), one set of output formats, and one last-run cache. %% lbl:HS-BOUND-04
  \item \texttt{HS-BOUND-05} --- Do not hardcode a config path and do not do filesystem I/O at import time. Keep pure constants at module top, defer resolution to an lru\_cache(maxsize=1) get\_settings() returning a frozen dataclass, try ordered strategies (env override, deployment descriptor walk-up, instance heuristic, system path) appending one audit entry each, warn and fall through on a stale override rather than crashing, and expose compat aliases through PEP 562 \_\_getattr\_\_. %% lbl:HS-BOUND-05
  \item \texttt{HS-BOUND-06} --- Do not give a tool family ad-hoc names or wire each tool in as its own MCP server. Name each tool <role>Upper with a binary short formed from the shortest stem prefix that stays unambiguous plus `up`, mirror it in the package as <short>\_pkg, and front the whole family with one router that discovers, verifies once, and fails closed on staleness. %% lbl:HS-BOUND-06
  \item \texttt{HS-IDENT-01} --- Do not ship a fast non-cryptographic hash as a key without a side-table, and do not reach for a cryptographic hash to avoid one. Bind identifier → name in a registry at registration, make re-registering the same name idempotent, reject a second name on an existing identifier loudly and name BOTH strings in the error, keep the string comparison off the hot path, and never reuse the identifier as an integrity or access primitive. %% lbl:HS-IDENT-01
  \item \texttt{HS-IDENT-02} --- Do not accept a validly signed descriptor that presents a new key for an existing name, and do not let a fleet share an unjittered TTL. Canonicalize bytes before signing or verifying, pin the signer fingerprint on first contact, reject a mismatch even when the signature verifies, never pin an unverified descriptor, make re-keying an operator action, and add uniform jitter to every refresh. %% lbl:HS-IDENT-02
  \item \texttt{HS-IDENT-03} --- Do not decide solution reuse on text similarity, and do not substitute a cached answer silently. Compile the problem to a typed graph with pre/postconditions, key it on an order-independent fingerprint, shortlist by similarity, break ties with a BOUNDED edit distance over the shortlist only, use separate shortlist and commit thresholds, score on several weighted components, emit an evidence record including the node correspondence, and gate store admission on novelty and validity. %% lbl:HS-IDENT-03
  \item \texttt{HS-EVID-01} --- Do not express a policy denial as a free-form string and do not rely on hand-picked cases to prove a refactor changed nothing. Map every rejection to a closed DenyReason enum, freeze the hook chain as an immutable versioned tuple evaluated first-deny-wins, record a large golden decision corpus, and fail CI on any divergence from it. %% lbl:HS-EVID-01
  \item \texttt{HS-EVID-02} --- Do not clear a quality gate by disabling the rule that found the defect. Baseline one measured count per defect class, recompute it live and isolated from project-local config, fail any commit whose count exceeds its baseline, record improvement without requiring it, and report the delta rather than the word 'regressed'. %% lbl:HS-EVID-02
  \item \texttt{HS-EVID-03} --- Do not adopt an optimization because it is implemented or because a paper reports a number. Open a ledger row per technique with its predicted uplift and quality floor, freeze a baseline on a named hardware fixture, implement behind an off-by-default flag, and set the verdict from artifacts: ADOPT or KILL, never blank — a KILL parks the code with a post-mortem, and every adopted technique is re-measured in combination. %% lbl:HS-EVID-03
  \item \texttt{HS-EVID-04} --- Do not execute a brief without auditing it, and do not state behaviour as prose. Open with a defect register giving each inherited defect an id and a specific mitigation, state every premise as a numbered claim with a citation and what the plan does with it, write conditions as logic blocks with a legend, name safety and liveness properties separately, bind each to a test id and each failure-mode property to a chaos drill that must be RUN, state phase exits as conjunctions, and put status in the filename verified against the code. %% lbl:HS-EVID-04
  \item \texttt{HS-LIVE-01} --- Do not decide liveness on a single binary timeout, and do not call the clock from inside the logic. Accrue misses across NORMAL ▸ SUSPECT ▸ DOWN so load can be shed before an outage is declared, reset and emit recovery immediately on any valid beat, inject the clock as a parameter defaulting to time.monotonic so the transition table is testable without sleeping, and never let a transition callback's exception stop the scheduler. %% lbl:HS-LIVE-01
  \item \texttt{HS-LIVE-02} --- Do not add a reaper, a reconnect protocol, or a persisted session registry to manage orphaned remote sessions. Bind the remote resource's lifetime to the control connection so every disconnection cause takes the same hangup path, keep session state in-process only, give teardown its own CLOSING state, never resume by identifier, write the rationale beside the non-persistence, and serve long-running work with a separate observer surface. %% lbl:HS-LIVE-02
  \item \texttt{HS-EXEC-01} --- Do not serve I/O, compute, and blocking work from one pool, and do not size or place workers from a core count. Import real machine topology and pin compute workers to cores, split pools by blocking character, size compute below the core count, state the yield budget and the re-acquire cadence as contract constants, put a task that cannot reach a checkpoint in the blocking pool rather than weakening the budget, bound the blocking pool, deliver completions through the data-path rings, and scrub child environments on inherit AND on exit. %% lbl:HS-EXEC-01
  \item \texttt{HS-EXEC-02} --- Do not poll or sleep to learn that native or accelerator work finished. Allocate an eventfd with EFD\_NONBLOCK|EFD\_CLOEXEC, register it with loop.add\_reader, release the interpreter lock before dispatching, have the native completion hook write the 8-byte counter, read exactly 8 bytes on wake and then drain a separate completion queue, and remove the reader before closing the descriptor. %% lbl:HS-EXEC-02
  \item \texttt{HS-EXEC-03} --- Do not compute variance as Σx² − (Σx)²/n, and do not recompute a sliding window from stored samples. Accumulate with Welford's recurrence, hold the window as two stacks each carrying its own aggregate, merge them with the exact two-group combination formula, pre-allocate both stacks so the ingest path allocates nothing, and reset accumulators when the window or sample clock is reconfigured. %% lbl:HS-EXEC-03
  \item \texttt{HS-SEC-01} --- Do not write plaintext secrets to persistent disk, and do not scrub logs by matching against secret VALUES. Decrypt into volatile memory, hold values in a read-only mapping whose repr names only keys, redact egress using the key set alone, render to an asserted tmpfs at 0600 with immediate unlink only when a consumer demands a path, AST-check both FunctionDef and AsyncFunctionDef for forbidden paths and writers, merge as `secret\_env | caller\_env` so an explicit caller value wins, and terminate on a decryption failure rather than proceeding with partial config. %% lbl:HS-SEC-01
  \item \texttt{HS-SEC-02} --- Do not rename foreign attributes, literals, or external keyword names when obfuscating a package, and do not drive the build from a shell script. Compute one global salted deterministic mapping over the package's BINDINGS only, subtract an explicit keep set that IS the published API, rewrite `import x` to `import x as \_alias`, drive the four stages from a real program that asserts no kept symbol was renamed before compiling and no source shipped after packaging, keep the package initializer as real source, build in an isolated environment, and never place a secret in the artifact. %% lbl:HS-SEC-02
  \item \texttt{HS-KNOW-01} --- Do not answer a structural code question from vector similarity alone, nor a natural-language one from graph traversal alone. Dual-write chunks to the vector store and AST relations to the property graph in one sync stage, resolve queries by taking vector anchors first and expanding along CALLS/INHERITS/IMPORTS with a hard hop cap, and de-duplicate identical source spans before assembling the prompt. %% lbl:HS-KNOW-01
\end{itemize}
% >>>END GENERATED<<<

\section{Invariants and Derivations (hand-curated)}

Every rule in the register above carries a \texttt{constraint} (what to do) and,
where formal content exists, an \texttt{invariant} (what is checkable). This
section derives the invariants. Each is anchored to its SSOT rule id, so a
change to a \texttt{[[rules]]} entry has exactly one place to land here.

\subsection*{Notation}

\begin{tabular}{@{}ll@{}}
\toprule
Symbol & Meaning \\
\midrule
$\Box P$        & $P$ holds at every point of every execution (safety) \\
$\Diamond P$    & $P$ holds at some point of the execution (liveness) \\
$\rho$          & fully symlink-resolving real-path function on paths \\
$\sqsubseteq$   & the safety order on actions: $a \sqsubseteq b$ iff $a$ is no more destructive than $b$ \\
$\bot$          & the unevaluable outcome (timeout, crash, absent evidence) \\
$K_3$           & Kleene three-valued logic over $\{\mathrm{true}, \mathrm{false}, \bot\}$ \\
$\varepsilon$   & unit roundoff of the floating-point format \\
$\mathrm{TV}(x)$& total variation of signal $x$ over the window under discussion \\
\bottomrule
\end{tabular}

\medskip
Symbols not listed are defined at first use. Where a rule has no formal content
beyond its constraint --- a naming convention, a directory layout --- this is
stated rather than dressed in notation it does not earn.

% ─────────────────────────────────────────────────────────────────────────
\section{GATE --- Admission and decision}

\subsection{Fail-closed totalisation}
\label{lbl:HS-GATE-01}

Let $D = \{\textsc{deny}, \textsc{allow}\}$ ordered $\textsc{deny} \sqsubseteq
\textsc{allow}$, and let a check be a partial function $c : \mathrm{Ctx}
\rightharpoonup D$, undefined exactly where it could not be evaluated. A gate is
the conjunction $g = \bigwedge_{i} c_i$. Totalising $g$ means choosing
$\gamma(\bot)$.

\begin{prop}
Fail-closed totalisation, $\gamma(\bot) = \textnormal{\textsc{deny}}$, is preserved by
conjunction and destroyed by disjunction.
\end{prop}

\begin{proof}
In $K_3$, $\bot \wedge \mathrm{false} = \mathrm{false}$ and $\bot \wedge
\mathrm{true} = \bot \mapsto \mathrm{false}$; either way the conjunction of a
fail-closed gate with anything is fail-closed. For disjunction, $\bot \vee
\mathrm{false} = \bot \mapsto \mathrm{false}$ but the \emph{intent} of a
disjunctive fallback is that an unevaluable primary defers to the alternative,
which reintroduces $\gamma(\bot) = \textsc{allow}$ at the point where the
alternative is itself unevaluable.
\end{proof}

The operational consequence is the invariant:
\begin{equation}
g(\mathrm{ctx}) = \textsc{allow} \;\longrightarrow\; \forall c \in
\mathrm{checks},\ \mathrm{evaluated}(c) \wedge \mathrm{passed}(c).
\end{equation}
An \textsc{allow} is therefore a positive claim about every check, which is why
the audit record must be written on accepts as well as rejects: the accept is
the assertion that needs evidence.

\subsection{Ordered negation yields a partition}
\label{lbl:HS-GATE-02}

Let admission be $S(x) \equiv \bigwedge_{i=1}^{n} c_i(x)$ with the conjuncts
declared in a fixed order. Then $\neg S = \bigvee_i \neg c_i$, which is
\emph{not} a partition of the rejected set: several conjuncts may fail on the
same $x$.

\begin{prop}
Define the reason function $r(x) = \min\{\, i : \neg c_i(x) \,\}$. Then
$\{\, r^{-1}(i) \,\}_{i=1}^{n}$ is a partition of $\neg S$, and consequently
$\sum_i \mathrm{count}_i = \mathrm{count}(\text{dropped})$.
\end{prop}

\begin{proof}
Exhaustive: $x \in \neg S$ implies the set $\{i : \neg c_i(x)\}$ is non-empty,
so its minimum exists and $x \in r^{-1}(r(x))$. Mutually exclusive: the minimum
of a non-empty subset of a totally ordered finite set is unique, so $r$ is a
function and the preimages are disjoint. Counters therefore neither
double-count nor lose a drop.
\end{proof}

Two corollaries drive the rules. First, \emph{the declared order is part of the
specification}: without it $r$ is not well defined and the counters are not
reproducible across runs. Second, every conjunct must be materialised in the
implementation even where it cannot fail locally, because the partition is
indexed by position --- a silently omitted conjunct shifts every later index.

\paragraph{Site disjointness.} Let $V_s$ and $V_w$ be the variables observable
at the sending and receiving sites. A gate is well-posed at a site only if
its conjuncts range over that site's variables:
\begin{equation}
\mathrm{Vars}(S) \subseteq V_s, \qquad \mathrm{Vars}(W) \subseteq V_w .
\end{equation}
A conjunct over remote state (``destination ring writable'') violates this and
is unimplementable without a coupling channel; the fix is to split one gate
into two over disjoint variable sets, not to add the channel.

\subsection{Delivery as a total function}
\label{lbl:HS-GATE-03}

The QoS classes make delivery a total function
\begin{equation}
\mathrm{deliver} : \mathrm{Class} \times \mathrm{RingState} \to \mathrm{Action},
\end{equation}
whose graph is exactly the truth table. Totality is the point: a library with
one global policy is the constant function in its first argument, which cannot
be correct for topics with different loss tolerances. Parameterising the test
over the table means the specification and the test are the same object, so a
new class adds rows to both or to neither.

The default is the $\sqsubseteq$-safest class. \textsc{lossless} pays latency;
\textsc{lossy} pays silence. Latency is observable by the party that suffers it,
silence is not, so the safe default is the one whose cost is visible. Note the
liveness obligation this creates: a \textsc{lossless} producer blocked on a
dead subscriber is a deadlock unless the block is bounded by a timeout or a
liveness signal from~\ref{lbl:HS-LIVE-01}.

\subsection{The hysteresis band bounds transition count}
\label{lbl:HS-GATE-04}

A single-threshold gate is memoryless, $s' = [\,x > \theta\,]$, so its output
transitions once per level crossing of $x$ at $\theta$. For $x = \mu +
\epsilon$ with $\mu \approx \theta$, the crossing rate is that of the noise
$\epsilon$ and grows without bound as the noise bandwidth grows. The two
threshold map is stateful:
\begin{equation}
s' = \begin{cases}
1 & x > \theta_{hi} \\
0 & x < \theta_{lo} \\
s & \theta_{lo} \le x \le \theta_{hi}.
\end{cases}
\end{equation}

\begin{prop}[Chatter bound]
Let $w = \theta_{hi} - \theta_{lo} > 0$. Over any window, the number of output
transitions satisfies
\begin{equation}
\#\mathrm{transitions} \;\le\; \left\lceil \frac{\mathrm{TV}(x)}{w} \right\rceil .
\end{equation}
\end{prop}

\begin{proof}
Consecutive output transitions are of opposite sense: after a transition to $1$
the next transition must be to $0$, which requires $x$ to fall from above
$\theta_{hi}$ to below $\theta_{lo}$, a traversal of variation at least $w$.
Each transition therefore consumes at least $w$ of the signal's total variation,
and the count is bounded by $\mathrm{TV}(x)/w$.
\end{proof}

The single-threshold gate is the limit $w \to 0$, where the bound diverges ---
which is precisely the observed flapping. Two design rules follow directly. The
band must be set from the measured noise amplitude, since $w$ appears in the
denominator and an unmeasured $w$ gives an unknown bound. And $\theta_{lo} <
\theta_{hi}$ must be enforced at construction, because equality silently
restores the divergent case.

Debouncing and input smoothing are rejected because they act on the wrong term:
debouncing delays every genuine transition by the debounce interval without
changing $\mathrm{TV}(x)$, and smoothing reduces $\mathrm{TV}(x)$ by attenuating
exactly the transients the detector exists to find.

\subsection{Consent as a conjunction with non-interfering denial}
\label{lbl:HS-GATE-05}

Export is the $K_3$ conjunction
\begin{equation}
\textsc{export} \equiv \mathrm{CONSENT} \wedge \mathrm{SCRUB\_OK} \wedge
\mathrm{NS\_POLICY}, \qquad \bot \mapsto \mathrm{false},
\end{equation}
so a consent service that times out denies rather than deferring to a cached
answer. The $K_3$ reading is stronger than short-circuit evaluation: $\bot
\wedge \mathrm{false} = \mathrm{false}$, so a definite failure elsewhere still
yields a definite denial and the counter attributes it correctly.

\begin{prop}[Denial non-interference]
Let $a$ range over artifacts and let $\mathrm{obs}_{req}$ be what the requester
observes. The consent-class denial is safe only if
\begin{equation}
\forall a_1, a_2 \quad \mathrm{obs}_{req}(\mathrm{reject}(a_1)) =
\mathrm{obs}_{req}(\mathrm{reject}(a_2)).
\end{equation}
\end{prop}

That is, the requester's view of a denial must be a constant function of the
artifact. A reason string is by construction a non-constant function of the
thing being protected, so \emph{the explanation is the disclosure}. The privacy
class is different in kind: there the owner already holds the data, so detail
flows to the owner and not to the requester. Both classes increment their own
counter, which preserves observability of the \emph{rate} without disclosing
any individual decision --- an aggregate over a protected predicate, which is
the standard construction for making a secret measurable.

Scan-before-persist follows from a containment argument rather than a logical
one: if the scan runs at export, the store holds unscanned content for the
entire interval between write and export, and the store is the artifact most
likely to be backed up and replicated. Moving the scan to the write path makes
$\mathrm{SCRUB\_OK}$ an invariant of the store rather than a property of one
egress path.

\subsection{Verdicts carry the repair function's argument}
\label{lbl:HS-GATE-06}

An exit code is a map onto $\{0,1\}$, carrying at most one bit. The repair an
agent must perform is a function $\mathrm{fix}(\mathrm{rule\_id},
\mathrm{file}, \mathrm{line}, \mathrm{measured}, \mathrm{expected})$. A function
cannot be recovered from a projection that discards its arguments, so an agent
handed only the exit code must guess --- and the cheapest guess that clears the
gate is to disable the rule. The structured verdict restores the arguments;
that is the whole mechanism, and it is why suppressing a finding is a doctrine
violation rather than a style preference.

Termination is a separate obligation. The repair relation is not known to be
contractive: an agent may re-emit a fix that fails the same rule. Termination
is therefore \emph{imposed} by a hard cap of $k$ re-runs per finding set,
followed by escalation. A cap is not a heuristic here; without it the loop has
no termination argument at all.

\subsection{Trust is a time-indexed predicate}
\label{lbl:HS-GATE-07}

Verification establishes $\mathrm{verified}(b, t_0)$ and records a fingerprint
$h(b)$ at $t_0$. The naive cache asserts $\mathrm{verified}(b, t)$ for all
later $t$, which is unsound because $b$ can change. Re-checking the fingerprint
before each call restores the invariant
\begin{equation}
\Box\big(\mathrm{call}(b) \rightarrow h(b) = h_{\mathrm{verified}}(b)\big),
\end{equation}
sound under collision resistance of $h$ against the adversary's budget. The
cost argument is that the expensive handshake is paid once per distinct
$(b, h(b))$ pair while the cheap fingerprint comparison is paid per call, so
verification is amortised to $O(1)$ per call without weakening the invariant.
This is why the row must carry the fingerprint and not merely the boolean:
staleness is undetectable from ``verified'' alone.

\subsection{Substitute-green is a necessary, not sufficient, condition}
\label{lbl:HS-GATE-08}

Let $I$ be the driver interface, $R$ the real backend, $M$ the substitute, and
$C$ a conformance suite: a set of predicates over $I$. Running $C$ against $M$
establishes $M \models C$. It does \emph{not} establish $R \models C$, and it
establishes $M \equiv R$ only up to $C$-expressible behaviour:
\begin{equation}
M \models C \ \wedge\ R \models C \;\Longrightarrow\; M \equiv_C R,
\end{equation}
where $\equiv_C$ is indistinguishability by the properties $C$ expresses.
Throughput and latency are not among them when $C$ runs against $M$, because
$M$'s timings are synthetic. Hence the two-part rule: substitute-green
$\models$ ordering and completion-count invariants, and substitute-green
$\nvDash$ throughput. Real-hardware ratification is a distinct later gate, and
a benchmark taken in substitute mode is not a hardware number --- which is why
the capability matrix must be \emph{printed}, so the mode is never inferred
from the logs.

The compile-time gate (default \textsc{off}) is preferred over a runtime import
guard because it removes the absent toolchain from the build's dependency
graph entirely, rather than making it a dependency that happens to be caught.

% ─────────────────────────────────────────────────────────────────────────
\section{MUT --- Persistent mutation and publication}

\subsection{Containment and omission monotonicity}
\label{lbl:HS-MUT-01}

Write $\rho$ for the fully symlink-resolving real-path function and
$\mathrm{Desc}(r)$ for the set of paths that are descendants of $r$ under
path-component decomposition. The security property is
\begin{equation}
\mathrm{contained}(c, r) \equiv \rho(c) \in \mathrm{Desc}(\rho(r)).
\end{equation}

\begin{prop}
The string-prefix test $\mathrm{str}(c) \succ \mathrm{str}(r)$ and
$\mathrm{contained}(c,r)$ are logically independent: neither implies the other.
\end{prop}

\begin{proof}
Prefix without containment, two ways. Non-alignment: take $c$ to be
\texttt{/var/log-backup} and $r$ to be \texttt{/var/log}. The string prefix
holds, but \texttt{log-backup} is not the component \texttt{log}, so
$c \notin \mathrm{Desc}(r)$. Symlink escape: $c = \texttt{/var/log/x}$ with $x$ a symlink
to \texttt{/etc} satisfies the string prefix while $\rho(c) = \texttt{/etc}
\notin \mathrm{Desc}(\rho(r))$. Containment without prefix: let $r$ itself be
reached through a symlink or bind mount, so $\rho(r)$ differs textually from
$r$ while $\rho(c)$ is a genuine descendant of $\rho(r)$.
\end{proof}

Because the two predicates are independent, no amount of hardening the string
comparison converts it into the security property. Only $\rho$ then the
descendant relation does, and an escape must be \emph{refused} rather than
followed --- following is exactly how a rooted tool acts outside its root.

\begin{prop}[Omission monotonicity]
A destructive tool's flag set is safe only if dropping any flag is
non-increasing in destructiveness:
\begin{equation}
\forall f \in F, \quad \mathrm{destructiveness}(\mathrm{inv} \setminus \{f\})
\ \sqsubseteq\ \mathrm{destructiveness}(\mathrm{inv}).
\end{equation}
\end{prop}

A \texttt{--force} flag satisfies this, since removing it moves the invocation
to simulation. A \texttt{--dry-run} flag violates it, since removing it moves
the invocation to deletion. The rule ``simulate by default'' is therefore not a
preference but the unique flag polarity under which the mode reached by
forgetting something is the safe one. The same argument condemns an inherited
configuration that can set \texttt{dry\_run\_default = false}: a safety default
a migration can switch off is not a default, because the migration is an
omission from the operator's point of view.

The remaining rules are corollaries of blast-radius separation. Discovery,
adoption and rotation are separate verbs so that the set of things that could
be deleted is reviewable before anything is; the report distinguishes
\texttt{deleted}, \texttt{would\_delete}, \texttt{refused} and \texttt{error} so
that a safety refusal is visible rather than collapsed into an error count; and
each step carries a timeout, because a cleanup pass hung on a stale mount
blocks the schedule and the next run overlaps it.

\subsection{Journal-before-effect is the only recoverable ordering}
\label{lbl:HS-MUT-02}

Let a multi-step mutation be a sequence of steps $s_1, \dots, s_n$. Write $A$
for the set of steps whose effect has been applied to the world and $J$ for the
set whose journal entry is durably committed. A crash may occur at any point.

\begin{prop}
Under the ordering \emph{commit then effect}, $A \subseteq J$ at every crash
point, and recovery is a total function of $J$. Under the reverse ordering,
$A \setminus J$ can be non-empty and recovery has no information about it.
\end{prop}

\begin{proof}
Commit-then-effect: an effect is applied only after its entry is committed, so
membership in $A$ implies membership in $J$; the crash truncates the trace, and
truncation preserves the implication. Since $J$ is durable and $A \subseteq J$,
the recovery procedure reading $J$ can decide for each entry whether to resume
or revert. Effect-then-commit: a crash in the window between the effect and the
commit places a step in $A \setminus J$; $J$ contains no record of it, and the
world does not distinguish ``step applied'' from ``step applied and later
overwritten'', so no function of $J$ recovers it.
\end{proof}

This is the entire safety argument, stated as the temporal property
\begin{equation}
\Box\big((\mathrm{effect\_applied}(s) \rightarrow \mathrm{entry\_committed}(s))
\wedge (\mathrm{entry\_committed}(s) \rightarrow \mathrm{replayable}(s))\big).
\end{equation}

Three consequences.

\paragraph{The inverse is recorded, not derived.} The inverse of a step is a
function of the state observed \emph{before} the step. At revert time that
state is gone, so an inverse derived then is a different function. It must be
computed while the pre-state is observable and stored in the entry.

\paragraph{Recovery is idempotent.} Recovery is itself a mutation and can be
interrupted, so the recovery function must satisfy $\mathrm{rec} \circ
\mathrm{rec} = \mathrm{rec}$; otherwise a crash during recovery is
unrecoverable by the same argument applied one level up.

\paragraph{Inverses are leaves.} If the inverse of an operation is itself
journalled as an invertible operation, the ``revert a revert'' relation is not
well-founded and revert recursion is unbounded. Classifying inverses as leaves
makes the relation well-founded and guarantees termination.

\paragraph{Chain integrity, and its limit.} With $h_i = H(h_{i-1}, e_i)$, any
modification of entry $j$ breaks the link at $j$ and, since the check is
sequential, localises the first break. The honest limit: truncation of a
\emph{suffix} is not detectable from the chain alone, because a shorter valid
chain is still a valid chain. Detecting suffix loss requires an external anchor
(a committed head pointer or an external witness), and a chain break must be
raised as an operational alert rather than parsed as a format error.

\subsection{Per-operation pinning is snapshot isolation}
\label{lbl:HS-MUT-03}

Let $S_0, S_1, \dots$ be the published snapshots and let a batch be $B = (e_1,
\dots, e_n)$.

\begin{prop}
Reading the current-snapshot pointer once per element yields only
\begin{equation}
\forall i\ \exists k \quad \mathrm{result}_i = f(e_i, S_k),
\end{equation}
whereas pinning once per operation yields the strictly stronger
\begin{equation}
\exists k\ \forall i \quad \mathrm{result}_i = f(e_i, S_k).
\end{equation}
\end{prop}

The quantifier order is the whole difference. Under the weaker form every
element is individually consistent with \emph{some} snapshot and the batch as a
whole corresponds to none, so the output is plausible and unattributable. This
is why the pin is acquired once and held for the operation's whole duration,
and why each snapshot carries an id recorded with the results: the question
``which version produced this?'' must have an answer.

Immutability buys the build algorithm as well as the absence of a lock. A
structure never mutated after publication can be bulk-loaded --- STR packing
for an R-tree --- which is unavailable to a structure that must accept
incremental insertion.

\paragraph{Reclamation.} Publication is one atomic pointer store, so readers
observe the old or new snapshot and never a mixture; multi-step publication
would have an observable intermediate. Freeing must satisfy
\begin{equation}
\mathrm{free}(S_k) \ \text{requires}\ |\mathrm{readers}(S_k)| = 0,
\end{equation}
which reference counting or epoch-based reclamation provides. Freeing at swap
time is a use-after-free whose trigger is load-dependent and therefore rare in
testing and common in production. The characteristic \emph{opposite} failure is
a handle leaked on an error path, retaining snapshots forever --- which is why
the chaos drill pairs mid-batch swapping with a flat-resident-memory soak.

\subsection{Head-advance-on-commit gives crash atomicity}
\label{lbl:HS-MUT-04}

The ring's observable state is a function of the committed lengths alone. Since
\texttt{reserve} does not advance the head and only \texttt{commit} does,
\begin{equation}
\text{failure anywhere in } (\mathrm{reserve}, \mathrm{commit}) \;\Longrightarrow\;
\mathrm{state}_{\mathrm{after}} = \mathrm{state}_{\mathrm{before}} .
\end{equation}
An exception between the two is a no-op, and a forgotten commit defaults to
$\mathrm{commit}(0)$, which is the identity on the head. The commit closure is
guarded by a \texttt{committed} flag, making commit idempotent, so the
context manager's default cannot double-advance a head the caller already
advanced.

Length is caller-provided because the ring is a slot allocator, not a stream
parser: it has no framing from which to infer where the payload ends. Inferring
it would require either a delimiter (which the caller would have to write
anyway) or a scan (which is the copy the pattern exists to avoid).

\paragraph{The locking discipline is the performance claim.} The critical
section is $\{\mathrm{reserve}, \mathrm{commit}\}$ --- $O(1)$ head arithmetic
--- while the caller's copy happens outside it. Serialised work per message
therefore drops from $O(|\mathrm{payload}|)$ to $O(1)$, which is what makes the
ring concurrent as well as zero-copy. A lock held across the copy would
serialise write bandwidth and forfeit most of the benefit.

\paragraph{Corruption is quarantined, not retried.} A ring that has failed
structurally is added to a process-local corrupted set and subsequently
no-ops via the health check. Retrying a structurally corrupt ring cannot
succeed; recovery is re-provisioning, which is an operator action.

\subsection{Epochs resolve the ABA ambiguity}
\label{lbl:HS-MUT-05}

A slot address alone is ambiguous across reuse: the pair $(\mathrm{ptr},
\mathrm{len})$ observed before and after a recycle are equal while denoting
different data. This is the ABA problem. Adjoining a monotone counter
$e$ --- a logical clock on the slot --- makes $(\mathrm{ptr}, e)$ unique over
the process lifetime, and the guard's check $e_{\mathrm{held}} =
e_{\mathrm{current}}$ is a happens-before test evaluated at every dereference.
A mismatch is a use-after-free hazard and must fail closed; returning the data
anyway is undefined behaviour dressed as a successful read. Views are read-only
because a writable view would let a consumer mutate a slot the producer has
already reclaimed.

\begin{prop}[Tail padding width]
A vector load of width $w$ beginning at any offset within an $n$-byte payload
may touch bytes up to index $n + w - 2$. Guaranteeing that every such access
lands in mapped memory therefore requires slot capacity
\begin{equation}
\mathrm{capacity} \;\ge\; n + (w - 1).
\end{equation}
\end{prop}

For AVX-512, $w = 64$, so 64 bytes of zeroed tail padding covers it with a byte
to spare and covers AVX2 ($w=32$) and NEON ($w=16$) a fortiori. Without the
padding the parser must either branch on the tail --- reintroducing the
bounds check the vectorisation exists to remove --- or risk a fault when the
payload ends near a page boundary.

% ─────────────────────────────────────────────────────────────────────────
\section{BOUND --- Interface boundaries}

\subsection{Types can express absence; they cannot express a postcondition}
\label{lbl:HS-BOUND-01}

Consider a sealing backend with the specification
\begin{equation}
\forall p \quad \mathrm{unseal}(\mathrm{seal}(p)) = p \ \wedge\ \mathrm{seal}(p) \neq p .
\end{equation}
The identity stub $\mathrm{seal}(p) = p$ satisfies the first conjunct and the
declared signature $(\mathtt{bytes}, \mathtt{str}) \to \mathtt{bytes}$, and
violates the second. The type system cannot express the second conjunct, so it
issues no signal: the stub type-checks and is catastrophically wrong.

The remedy is to move the information the type system \emph{can} carry. Let the
public return type be the sum
\begin{equation}
\mathrm{Resolved} = \mathrm{Present}(\mathrm{backend}) \;+\;
\mathrm{Absent}(\mathrm{capability}, \mathrm{reason}).
\end{equation}
Because $\mathrm{Absent}$ has no member \texttt{seal}, every call site that
uses the capability without discriminating is ill-typed. The check moves from a
reviewer to the compiler.

Three refinements follow. $\mathrm{Absent}$ must be falsy but distinct from
\texttt{None}, because \texttt{None} is the shared inhabitant of every optional
in the program: unifying with it collapses the coproduct tag and discards the
reason. The reason must be carried, since ``not available'' without a cause
becomes a support ticket. And the union must be the \emph{public} return type
--- narrowing it inside the library and raising on absence relocates the policy
decision away from the caller who owns it.

Typed degradation is still silent to an operator, so the resolved capability set
must be printed in the same doctor surface that reports versions and health.

\subsection{The facade is definitional equality}
\label{lbl:HS-BOUND-02}

A characterisation suite $C$ recorded against $M_{\mathrm{old}}$ in place, then
replayed against $M_{\mathrm{new}}$, establishes agreement on the behaviours
$C$ exercises --- and nothing more. The strength of the behaviour-preservation
argument is therefore exactly the coverage of $C$, which is why $C$ must be
recorded \emph{before} the move (a suite written afterwards records the new
behaviour) and why it must record behaviour the author considers wrong rather
than asserting what the code should do. Correcting a bug is a separate commit
with its own test.

\begin{prop}
A facade defined as re-export, $\mathtt{old}.x \equiv \mathtt{new}.x$, admits
one implementation under two names, so drift between the paths is not merely
unlikely but unrepresentable. A facade containing logic admits two
implementations, and their equality becomes an obligation that must be
maintained.
\end{prop}

The deprecation window is stated in the shim itself because a shim without a
stated removal release has no termination condition and is therefore permanent.
Both paths run the same suite during the window, since testing only the new
path leaves the shim unverified for exactly as long as callers still depend on
it. The public surface is frozen as a version-controlled file asserted by a
snapshot test, so an addition or removal appears as a diff a reviewer reads
rather than as a silent widening.

\paragraph{Ordering.} Extraction order must be a topological order of the
module dependency DAG. Extracting a module before its dependencies produces a
facade that imports backwards, i.e.\ a cycle, and the application does not
remain working between steps.

\subsection{Structural contract, bounded discovery}
\label{lbl:HS-BOUND-03}

Plugin discovery has little formal content, and the register says so. The two
statements worth making precise are these. The host--plugin contract is a
\emph{structural} type on the root module --- \texttt{app} required, the
lifecycle hooks optional --- rather than a nominal one, which is what allows a
plugin to be built without importing a base class from the host and therefore
without a build-time dependency on it. And the per-plugin import timeout bounds
host startup at $O(n \cdot t_{\max})$ for $n$ plugins; without it a single hung
\texttt{\_\_init\_\_.py} stalls the host's CLI and API heads indefinitely,
making startup an unbounded function of third-party code.

The PEP 420 namespace is what lets several independently built wheels
contribute to one \texttt{plugins} namespace without any of them owning it.

\subsection{One core removes a quadratic obligation}
\label{lbl:HS-BOUND-04}

With $n$ independently implemented surfaces, each pair must be kept consistent
in configuration schema, precedence, and output format, giving
\begin{equation}
\binom{n}{2} = \frac{n(n-1)}{2}
\end{equation}
hand-maintained consistency obligations --- three for CLI, MCP and TUI. Routing
every surface through one core reduces this to zero by construction, since
there is one schema and one formatter to be consistent with.

Precedence must be a \emph{total} order on configuration sources (CLI arguments
$>$ project TOML $>$ user config $>$ built-in defaults). A partial order leaves
incomparable pairs, and resolution of an incomparable pair is by definition
unspecified --- which surfaces as a setting that depends on load order.

\subsection{Deferred resolution makes import total}
\label{lbl:HS-BOUND-05}

A module that performs filesystem I/O at import time makes \texttt{import} a
partial function of the filesystem: it succeeds or fails depending on state the
importer did not choose, which is why such modules are un-importable under test
and in a fresh deployment. Deferring all I/O to first call makes import total,
and $\mathtt{lru\_cache(maxsize=1)}$ makes the resolution idempotent and $O(1)$
after the first call. The frozen dataclass makes the resolved settings
immutable, so a later mutation cannot silently diverge from what was resolved.

The audit list is the non-obvious part. The diagnostic value of the failure is
a function of the \emph{whole} strategy sequence, not of the last failure: an
error reporting only ``config not found'' requires the operator to re-run with
instrumentation to learn which paths were tried and why each was rejected.
Appending exactly one entry per strategy --- including on the strategies that
were skipped --- makes the message self-sufficient. A stale explicit override
warns and falls through rather than raising, because the operator's intent is
recoverable from the later strategies while a hard failure discards it.

\subsection{The short name is a prefix code}
\label{lbl:HS-BOUND-06}

Let $\mathrm{stem}(x)$ be the tool's role word, so
$\mathrm{stem}(\texttt{routeUpper}) = \texttt{route}$. The naming rule maps
\begin{equation}
\mathrm{short}(x) = \mathrm{prefix}(\mathrm{stem}(x)) \,\|\, \texttt{"up"} .
\end{equation}

\begin{prop}
The single-letter form $\mathrm{prefix} = \mathrm{head}$ is \emph{not}
injective on the current roster.
\end{prop}

\begin{proof}
$\mathrm{head}(\texttt{test}) = \mathrm{head}(\texttt{tailscale}) = \texttt{t}$
and $\mathrm{head}(\texttt{vpn}) = \mathrm{head}(\texttt{validate}) =
\texttt{v}$. Under the single-letter rule both pairs would collide on
\texttt{tup} and \texttt{vup}.
\end{proof}

The roster resolves this by lengthening the prefix on the later claimant:
\texttt{tsup}, \texttt{vpnup}, \texttt{wup}. The operative rule is therefore
\emph{shortest distinguishing prefix}, which makes the short names a prefix
code over stems and restores injectivity. This is worth stating because the
narrative document states the single-letter rule, which its own table
contradicts; the register carries the corrected form.

Injectivity is the property that matters, since the short name is what a
person types and what the router resolves. The router itself is
\ref{lbl:HS-GATE-07}: one connection, discovery, verify-once, and a staleness
re-check before each proxied call.

% ─────────────────────────────────────────────────────────────────────────
\section{IDENT --- Identity, keys and signatures}

\subsection{Injectivity is enforced, not assumed}
\label{lbl:HS-IDENT-01}

A hash $h : \mathrm{Names} \to 2^{64}$ cannot be injective, since its domain is
unbounded and its codomain is not. The registry does not improve $h$; it
restricts the domain to the registered set and enforces injectivity there.

\begin{prop}
Let $R_k$ be the registered set after $k$ registrations, where a registration
of $b$ is admitted iff $h(b) \notin h(R_k)$, or $b \in R_k$ (idempotent
re-registration). Then $h|_{R_k}$ is injective for every $k$.
\end{prop}

\begin{proof}
By induction. $R_0 = \emptyset$, vacuously injective. Assume $h|_{R_k}$
injective. A re-registration leaves $R_k$ unchanged. A new $b$ is admitted only
when $h(b) \notin h(R_k)$, so for all $a \in R_k$, $h(a) \neq h(b)$, and with
the inductive hypothesis $h|_{R_k \cup \{b\}}$ is injective.
\end{proof}

The invariant is thus
\begin{equation}
\forall a, b \quad \big(\mathrm{registered}(a) \wedge \mathrm{registered}(b)
\wedge h(a) = h(b)\big) \rightarrow a = b,
\end{equation}
a property of the registry's admission rule and not a hope about the hash.

\paragraph{What the side-table buys, quantitatively.} By the birthday bound,
$m$ names over a 64-bit codomain collide with probability approximately
$m^2 / 2^{65}$; for $m = 10^4$ that is about $2.7 \times 10^{-11}$. The
registry does not reduce this probability. It changes the \emph{consequence}:
without the table, a collision is a silent cross-delivery between unrelated
topics --- simultaneously a correctness and a confidentiality failure --- and
with it, a collision is a loud rejection at bind time, when a person is present
to rename a topic. Converting a small probability of silent corruption into the
same small probability of a loud, actionable error is the entire value, and it
is why reaching for a cryptographic hash is the wrong move: it pays a per-
message cost on a path chosen for speed and still does not detect the collision
if one occurs.

The error must name both colliding strings, since the operator's remedy is to
rename one and they cannot choose without knowing which two. Re-registration of
the same name is idempotent, or process restarts and duplicate subscriptions
become spurious failures. And the identifier is a routing key: it must not be
reused as an integrity check, an access token, or a cross-boundary dedup key,
all three of which assume collision resistance the function does not provide
and the table supplies only within one registry's scope. Registry scope is
therefore a trust boundary --- two independently built registries may bind the
same identifier to different names.

\subsection{Signature validity is not authority to re-key}
\label{lbl:HS-IDENT-02}

The verifier's predicate $\mathrm{verify}(m, k)$ asserts that $m$ was signed by
the holder of $k$. It asserts nothing about whether $k$ is \emph{the} key for a
given name. Conflating the two is exactly the vulnerability: an adversary
holding any valid key can then rebind a name. Pinning separates them, adding
the binding $\mathrm{name} \mapsto \mathrm{fingerprint}(k)$ established at
first contact and enforced thereafter, so that
\begin{equation}
\mathrm{accept}(\mathrm{name}, m, k) \equiv \mathrm{verify}(m,k) \wedge
\big(\mathrm{fingerprint}(k) = \mathrm{pin}(\mathrm{name})\big).
\end{equation}

\paragraph{The trust model, stated honestly.} TOFU converts a requirement for
continuous authentication into a single point-in-time assumption: the scheme is
sound iff the adversary is absent at $t_0$. It \emph{detects} a later key change
and does not \emph{prevent} one, which is why re-keying must be an explicit
operator-driven pin reset rather than an automatic acceptance. An unverified
descriptor is never pinned, or the assumption at $t_0$ is not even established.

\paragraph{Canonicalisation.} Verification compares bytes, so signer and
verifier must agree on the byte representation exactly --- sorted keys, compact
separators, fixed escaping. Without canonicalisation, an equal-as-JSON message
is unequal as bytes and verification fails nondeterministically across
implementations.

\paragraph{Jitter.} $N$ nodes sharing a static TTL $T$ refresh at the times
$\{t_0 + kT\}$, so demand is a comb of impulses of height $N$. Adding
independent jitter $J \sim U[0, j]$ spreads each impulse over a window of width
$j$, giving expected instantaneous demand $N/j$. The peak load falls by roughly
$j/\Delta$ where $\Delta$ is the service window, so $j$ must be chosen
comfortably larger than the service time; a jitter smaller than the service
time does not decorrelate the herd.

\subsection{The fingerprint is a filter, the edit distance is the decision}
\label{lbl:HS-IDENT-03}

Weisfeiler--Lehman hashing is invariant under graph isomorphism, which is what
makes it usable as a key at all: without order-independence the same problem
compiled from differently ordered input misses its own cache entry, and the
cache degrades to useless while continuing to appear to work.

The converse fails, and the pipeline is built around that fact.

\begin{prop}
$\mathrm{WL}(G_1) = \mathrm{WL}(G_2)$ does not imply $G_1 \cong G_2$; the
1-dimensional WL refinement cannot distinguish, for instance, two
non-isomorphic regular graphs of the same order and degree.
\end{prop}

The fingerprint is therefore a \emph{necessary} condition for a match, not a
decision procedure. This is precisely why a bounded structural edit distance
runs downstream on the shortlist: the cheap complete-recall filter narrows the
field, and the expensive precise comparison decides. Running the edit distance
against the whole store would be correct and unaffordable; running the
fingerprint alone would be affordable and wrong.

\paragraph{Bounded cost.} With graph size pruned at compile time, shortlist
width $k$ fixed, and the edit-distance beam bounded, total query cost is
\begin{equation}
O(\mathrm{compile}) + O(\mathrm{ANN\ retrieve}) + k \cdot O(\mathrm{edit}),
\end{equation}
which is constant in $|\mathrm{store}|$ given sublinear approximate retrieval.
Every stage must be bounded; one unbounded stage makes worst-case latency grow
with the store, which is the failure that appears only after the store is large.

\paragraph{Composite score.} The score $s = \sum_i w_i c_i$ over fingerprint
similarity, inverse edit distance, precondition compatibility and I/O shape
agreement has two properties a single distance lacks: no one component can
carry a match alone, and when a reuse was wrong, $\arg\max_i w_i c_i$ identifies
the component that overweighted. Weights must therefore be visible. The two
thresholds --- lower for shortlist admission, higher for committing to reuse
--- are the band of \ref{lbl:HS-GATE-04} applied to a non-alerting decision,
and buy the same stability: a candidate near the boundary does not alternate
between reuse and recomputation across runs.

Evidence, including the node correspondence, is emitted on every match. The
correspondence is the part most often omitted and most valuable, since it is
what lets a person see which parts of the two problems the engine considered
equivalent. Admission gates on novelty and validity matter more than the
tiering: a store with no novelty gate fills with near-duplicates, and one with
no poison gate will faithfully reuse a bad solution forever.

% ─────────────────────────────────────────────────────────────────────────
\section{EVID --- Evidence and measurement}

\subsection{Replay bounds the disagreement rate}
\label{lbl:HS-EVID-01}

Let $f_{\mathrm{old}}, f_{\mathrm{new}} : \mathrm{Inputs} \to \mathrm{Verdicts}$
be the reference and candidate policy chains. A unit suite checks
$f_{\mathrm{new}}$ on a hand-picked set $T$ with $|T|$ small and chosen by the
same person who wrote the code. Differential replay checks
$f_{\mathrm{new}} = f_{\mathrm{old}}$ on a corpus $X$ drawn from the production
distribution $D$.

\begin{prop}[Rule of three]
Observing zero mismatches over $n$ independent samples from $D$ gives the
one-sided 95\,\% confidence bound
\begin{equation}
\Pr_{x \sim D}\big[f_{\mathrm{new}}(x) \neq f_{\mathrm{old}}(x)\big] \;\le\;
\frac{3}{n}.
\end{equation}
\end{prop}

\begin{proof}
If the true rate is $p$, the probability of observing no mismatch in $n$
independent trials is $(1-p)^n$. Setting $(1-p)^n = 0.05$ and using
$(1-p)^n \approx e^{-np}$ gives $np \approx -\ln 0.05 \approx 2.996$.
\end{proof}

For $n = 10^5$ the bound is $3 \times 10^{-5}$. This is a statement about
agreement with the previous behaviour on the production distribution, which is
what ``no silent drift'' means; it is not a statement that either function is
correct.

\paragraph{The closed taxonomy is a precondition of the gate, not a separate
nicety.} Comparison requires equality on verdicts. With free-form strings,
verdict equality degenerates to string equality on human-readable messages, so
rewording a message reports as drift. A gate that fires on cosmetic changes is
disabled within a quarter. Restricting rejections to a finite \texttt{DenyReason}
enum makes the codomain finite and comparable, and simultaneously lets a
downstream handler distinguish a policy denial from an internal error --- which
free-form strings do not.

\paragraph{First-deny-wins.} Evaluating the chain as a short-circuiting fold
makes the recorded reason the first failing hook, which is the same
minimum-index construction as \ref{lbl:HS-GATE-02} and yields the same
partition property. Freezing the chain as an immutable tuple tagged with a
monotonically increasing version is what makes ``the same chain'' a checkable
claim across replays.

\subsection{The ratchet is monotone, conditionally}
\label{lbl:HS-EVID-02}

Let $m_t$ be the measured count of a defect class at commit $t$ and $b_t$ its
baseline. The gate admits a commit iff $m_t \le b_t$, and on admission sets
$b_{t+1} = \min(b_t, m_t)$.

\begin{prop}
$(b_t)$ is non-increasing and bounded below by $0$, hence convergent; and no
admitted commit increases any metric.
\end{prop}

This is a ratchet in the mechanical sense: motion in one direction only. It is
strictly stronger than a fixed ceiling $c$, under which $m$ may drift freely
anywhere below $c$ and the accumulated debt is invisible until it crosses.

\paragraph{The soundness condition, which is easy to miss.} The ratchet
compares counts produced by a measurement \emph{procedure}. If the procedure is
allowed to read project-local configuration, then disabling a lint rule lowers
the count and the ratchet passes --- the exact failure it was built to prevent.
The ratchet is therefore sound only when the procedure is pinned and isolated:
an explicit rule selection, and \texttt{--isolated} / \texttt{--no-cache} so
that neither project config nor a stale cache can influence the number. Under
that condition, changing what is measured requires editing the baseline file,
which is a reviewable diff. Without it the ratchet measures the configuration
rather than the code.

One metric per defect class, because a metric conflating two classes cannot be
targeted for fix; and report the delta rather than the word ``regressed'',
because the delta is what identifies the change responsible.

\subsection{Composition is not the sum of isolates}
\label{lbl:HS-EVID-03}

Adoption is the conjunction
\begin{equation}
\textsc{adopt} \iff (\mathrm{uplift} \ge u_0) \wedge (\mathrm{quality} \ge q_0)
\wedge \neg\,\mathrm{composition\_regression},
\end{equation}
with the quality floor separate from the uplift floor precisely so that latency
cannot be bought with accuracy silently.

\begin{prop}
For $k$ candidate techniques there are $2^k$ enable-sets. Measuring the $k$
singletons leaves $2^k - k - 1$ configurations unmeasured, including the
full set --- which is the one that ships.
\end{prop}

Define the interaction term
\begin{equation}
\Delta \;=\; \mathrm{uplift}\Big(\bigcup_i \{i\}\Big) \;-\; \sum_i
\mathrm{uplift}(\{i\}).
\end{equation}
$\Delta = 0$ only if the techniques are independent. Optimisations competing for
one hot path typically share a bottleneck, making the effect subadditive
($\Delta < 0$): each removes part of the same cost, so the second cannot remove
it again. The composition matrix is therefore not diligence, it is the only
measurement of the shipped configuration.

\paragraph{Default-off is the gate.} A flag on by default before its row reads
\textsc{adopt} has already skipped the gate, and the shipped system is then a
configuration nobody measured.

\paragraph{Killed code is parked with a post-mortem.} Absence of code is not an
explanation, and without the record the same technique is re-proposed and
re-implemented by someone who only saw that it was missing.

\subsection{Prose has no negation operator}
\label{lbl:HS-EVID-04}

The mechanical case for logic blocks is one line:
\begin{equation}
\neg \bigwedge_i c_i \;=\; \bigvee_i \neg c_i .
\end{equation}
A gate written as an expression can be negated systematically, and the failure
taxonomy is \emph{derived} from it rather than invented alongside it --- which
is what keeps specification and error handling from drifting apart
(\ref{lbl:HS-GATE-02}). ``The service starts once its dependencies are
healthy'' admits at least three readings and supports no such derivation.

\begin{prop}[Alpern--Schneider]
Every trace property is the intersection of a safety property and a liveness
property.
\end{prop}

Naming safety $\Box P$ and liveness $\Diamond Q$ separately is therefore the
canonical decomposition rather than a stylistic preference, and it has a direct
testing consequence: a safety property is falsifiable by a finite execution
prefix, so a unit test can refute it; a liveness property is not falsifiable by
any finite prefix, since the good thing may still happen later. Liveness can
only be exercised by bounded-time observation under injected failure --- which
is exactly what a chaos drill is, and why the plan must state that the drill is
to be \emph{run} rather than merely specified.

The remaining rules are about provenance. Each premise is a numbered claim with
a citation \emph{and} a statement of what the plan does with it, so a reviewer
can check the inference and not only the source; an uncited number becomes a
target nobody can defend. The defect register makes the plan's disagreements
with its own mandate part of the permanent record, which a review comment
cannot. Phase exits are conjunctions of checkable conditions, because ``phase 2
complete'' is a claim while a conjunction is a gate. And status lives in the
filename and is verified against the code: in-document markers record what an
author believed when they typed them, the code records what is true now, and
when they disagree the code wins.

% ─────────────────────────────────────────────────────────────────────────
\section{LIVE --- Liveness and lifetime}

\subsection{Accrual trades false alarms against latency, exponentially}
\label{lbl:HS-LIVE-01}

A binary timeout is a single threshold on elapsed time and inherits the
chatter of \ref{lbl:HS-GATE-04}: under jitter or a garbage-collection pause the
signal crosses the boundary repeatedly, and each crossing is a failover
decision. The accrual detector replaces it with a monotone step function of the
miss count $M$, with thresholds $k_s < k_d$:
\begin{equation}
\mathrm{state}(M) = \begin{cases}
\textsc{down} & M \ge k_d \\
\textsc{suspect} & k_s \le M < k_d \\
\textsc{normal} & M < k_s .
\end{cases}
\end{equation}

\begin{prop}
If beats are lost independently with probability $p$, then the probability of
reaching a threshold $k$ with the peer alive is $p^{k}$, while the detection
latency at that threshold is $k \cdot \mathrm{interval}$.
\end{prop}

The false-alarm rate therefore falls \emph{exponentially} in the same parameter
that raises detection latency only \emph{linearly}, which is why a small
increase in $k$ is usually the right trade. Choose $k_d$ so that $p^{k_d}$ is
acceptable and accept $k_d \cdot \mathrm{interval}$ as the detection time; the
intermediate \textsc{suspect} state exists so load can be shed or a backup
warmed during the interval in which the evidence is suggestive but not
conclusive.

Recovery is deliberately asymmetric: any valid beat resets $M$ to zero and
emits a recovery transition immediately. Degradation requires accumulated
evidence, recovery requires one positive observation, because a beat is
conclusive proof of liveness while a miss is only weak evidence of death.

\paragraph{Clock injection restores referential transparency.} A class that
calls the clock internally is a function of ambient state, so its transition
table can be exercised only by real waiting --- slow, and flaky because the
wait must exceed the interval by an unknown margin. Injecting
$\mathrm{clock} : () \to \mathbb{R}$ makes time an input; the detector becomes
a pure function of (clock values, events) and the table is exhaustively
testable in $O(|\mathrm{table}|)$ evaluations with no sleeping. This is the
same discipline as the doctrine's monotonic-clock rule, and the default must be
\texttt{time.monotonic} since the quantity is an elapsed interval.

Finally, an exception raised by a transition callback must be caught and logged
without interrupting the scheduler: the detector's liveness must not depend on
the correctness of a consumer's handler.

\subsection{Orphans removed rather than reconciled}
\label{lbl:HS-LIVE-02}

Let $C$ be the controller's set of live sessions and $W$ the worker's. With
independent lifetimes the two diverge under any network event, and $W \setminus
C$ --- the orphans --- must be reconciled by a reaper (which guesses a timeout,
so it either kills live sessions or leaves dead ones), a reconnect protocol
(which must reconcile buffered output, cursor and process state across a gap in
which anything could have happened), or a persisted registry (which outlives
the processes it describes and becomes a list of identifiers pointing at
nothing). Each is a mechanism for \emph{managing} orphans.

Coupling the remote resource's lifetime to the control connection makes $W$ a
function of the live connection set, so
\begin{equation}
W \setminus C = \emptyset
\end{equation}
holds by construction. This is the ``make the invalid state unrepresentable''
move: the reconciliation problem is not solved, it is deleted.

\paragraph{Non-persistence is a lifetime argument.} Persisted state has lifetime
$\supseteq$ the process lifetime and can therefore outlive its referent;
in-process state has lifetime $=$ the process lifetime and cannot. Because
non-persistence looks like an omission, the rationale must be written beside it
or a later maintainer helpfully adds durability and reintroduces the orphan
class.

\paragraph{One path for every disconnection cause.} A deliberate close, a
crashed controller and a lost route all converge on the same hangup. One path
means the rare failure case is exercised by the common one, every day.

\paragraph{Teardown has duration.} A two-state model $\{\textsc{open},
\textsc{closed}\}$ asserts that hangup is instantaneous, which is false, and
the assertion becomes a race in which a caller opens a new session while the
old process still holds the terminal. The state machine needs one state per
distinguishable observable condition, hence \textsc{closing} alongside
\textsc{closed} and \textsc{errored}, with both terminal states marking the
cursor final so a reader can distinguish a clean end from a bad one without
either being resumable.

\paragraph{Refusing resumption removes an authorisation question.} A reconnect
protocol must decide whether a party presenting an old identifier is entitled
to the output buffered under it. Binding a UUID4 identifier to one connection's
lifetime and never re-attaching it removes the question rather than answering
it --- a security property obtained by removing a capability. The legitimate
need for long-running remote work is then met by a separate observer surface,
rather than by weakening the session contract for everyone.

% ─────────────────────────────────────────────────────────────────────────
\section{EXEC --- Execution substrate}

\subsection{One pool cannot satisfy three constraint sets}
\label{lbl:HS-EXEC-01}

Let $n$ be a pool's worker count. Compute work wants $n \approx$ the core count,
since additional workers only add context switching to work that is already
CPU-bound. Blocking work wants $n \gg$ the core count, since workers are
descheduled while waiting and a small pool serialises unrelated waits. These
constraints have empty intersection, so a single pool is misconfigured for at
least one class by construction. Splitting into an event-loop I/O pool, a
core-pinned compute pool and a bounded blocking pool makes the constraint set
satisfiable --- and the split is by \emph{blocking character}, not priority:
priority is a scheduling policy, while blocking character determines which pool
can survive the task at all.

\paragraph{Topology, not core count.} $\mathtt{cpu\_count()}$ is a cardinality.
Placement decisions depend on relations it does not encode: which cores are
hyperthread siblings of one physical core, which share an L3, and which NUMA
node owns the memory a worker's queue lives in. Importing the real topology and
pinning with an explicit affinity mask is what makes the queue's memory local
and stops two cache-hostile workers landing on siblings.

\paragraph{Sizing.} The compute pool is $N-1$, not $N$, so the event loop, the
blocking pool's supervisor and the operating system are not competing with a
pinned worker for the last core.

\paragraph{Work stealing.} With a Chase--Lev deque the owner pops from one end
and a thief steals from the other, so the common case --- a worker servicing
its own queue --- touches no contended cache line, and stealing costs only when
there is an imbalance to correct. Under the Blumofe--Leiserson analysis the
expected number of steals is $O(P \cdot T_\infty)$ for $P$ processors and
critical-path length $T_\infty$, i.e.\ independent of total work $T_1$: the
overhead scales with the shape of the computation, not its size.

\paragraph{Preemption is two explicit constants.} With yield budget $b$ and
re-acquire cadence $c$:
\begin{equation}
\text{worst-case co-scheduling delay} = b, \qquad
\text{worst-case signal latency} = c .
\end{equation}
Both are contract constants derived from a latency requirement, not
implementation details, which is why they are stated ($b = 5$\,ms, $c =
25$\,ms) rather than tuned silently. A checkpoint the task's own code must call
is a contract with the task's author; a task that cannot be chunked to reach a
checkpoint therefore belongs in the blocking pool, and the budget must not be
weakened to accommodate it.

\paragraph{Completions travel the data path.} A promise or future allocates on
the managed heap and synchronises with the runtime --- on precisely the path
the pool released the interpreter lock to avoid. Publishing into a preallocated
lock-free ring slot has neither cost, and the event loop already has a reader
registered there, so the wake-up path is the one used for data. The completion
mechanism must be in the same cost class as the work it completes, or it undoes
the lock release.

\paragraph{Bounding and scrubbing.} The blocking pool exists to quarantine
unbounded latency; unbounded, it quarantines nothing. And child environments
are scrubbed on inherit \emph{and} on exit, since scrubbing only the inherited
environment leaves the child's own exported values and its error output as an
egress path.

\subsection{Polling has a constant cost--latency product}
\label{lbl:HS-EXEC-02}

Let $T$ be a polling period over a wait of duration $d$. The number of wakeups
is $d/T$ and the expected notification latency is $T/2$, so
\begin{equation}
\mathrm{cost} = O(d/T), \qquad \mathrm{latency} = O(T), \qquad
\mathrm{cost} \times \mathrm{latency} = O(d) .
\end{equation}
The product is independent of $T$: there is no polling period that is both
cheap and prompt, only a choice of where on the hyperbola to sit. Readiness
notification breaks the trade --- cost is $O(\#\mathrm{events})$ and latency is
the kernel wake latency, independent of it. This is the reason the pattern
exists, and it is why ``just poll faster'' is not a fix.

\paragraph{eventfd semantics constrain the reader.} A non-semaphore eventfd is
a single 64-bit counter: concurrent writes \emph{add}, and a read returns the
accumulated value and resets it to zero. Writes therefore coalesce, and the
value read means ``at least one completion occurred since the last read'',
never ``exactly this many items are pending''. The correct structure is thus
the one the pattern prescribes: read the eight bytes to re-arm the descriptor,
then drain a separate completion queue for the actual work items. Treating the
counter as a work count loses completions whenever two writes land between
reads.

The flags are not optional. \texttt{EFD\_NONBLOCK} prevents the read from
blocking the event loop when a spurious wakeup finds the counter already
drained; \texttt{EFD\_CLOEXEC} prevents the descriptor leaking into children.
The reader must be unregistered before the descriptor is closed, or
\texttt{epoll} retains a reference to a closed fd. And the interpreter lock
must be released before dispatch, since the entire purpose is to let native
threads and accelerator streams run concurrently with the loop.

\subsection{Welford, and why the window needs two stacks}
\label{lbl:HS-EXEC-03}

\paragraph{Cancellation.} The textbook form computes
$\sum x_i^2 - (\sum x_i)^2 / n$. Both terms have magnitude $\approx n(\mu^2 +
s^2)$ while their difference has magnitude $\approx n s^2$, so the relative
error of the result is amplified by
\begin{equation}
\kappa \;=\; \frac{\mu^2 + s^2}{s^2} \;=\; 1 + \frac{\mu^2}{s^2}.
\end{equation}
For a sensor with $\mu / s = 10^3$ --- a small variance about a large mean,
which is the normal case in telemetry --- $\kappa \approx 10^6$, so about six
significant digits are lost before any accumulation error. Welford's recurrence
never forms that difference:
\begin{equation}
\delta = x - \bar{x}_{n-1}, \quad
\bar{x}_n = \bar{x}_{n-1} + \frac{\delta}{n}, \quad
M_2 \mathrel{+}= \delta \cdot (x - \bar{x}_n),
\end{equation}
accumulating $\sum (x_i - \bar{x})^2$ directly from quantities of the correct
magnitude.

\paragraph{Exact merge.} Two independent accumulators combine in constant time
(Chan, Golub and LeVeque):
\begin{equation}
n = n_a + n_b, \quad \delta = \mu_b - \mu_a, \quad
\mu = \mu_a + \delta\,\frac{n_b}{n}, \quad
M_2 = M_{2,a} + M_{2,b} + \delta^2\,\frac{n_a n_b}{n}.
\end{equation}
The correction term is the between-group contribution to the total sum of
squares --- the identity is the usual within/between decomposition of variance,
which is why the merge is exact rather than approximate.

\paragraph{Why two stacks, and not a ring buffer.} The obvious sliding-window
implementation subtracts the departing sample from the accumulator. That
requires inverting Welford's update, which is not numerically viable: the
inverse re-introduces a difference of nearly equal quantities and accumulates
error without bound over a long-running window, since nothing ever resets it.
The two-stack queue avoids inversion entirely --- aggregates are only ever
\emph{added} --- at the cost of a periodic rebuild.

\begin{prop}[Amortised $O(1)$]
Let $\Phi$ be the number of elements on the in-stack. Push costs $1$ and raises
$\Phi$ by $1$: amortised $2$. A pop with a non-empty out-stack costs $1$ and
leaves $\Phi$ unchanged: amortised $1$. A pop that first flips costs
$\Phi + 1$ and reduces $\Phi$ to $0$: amortised $1$. Every operation is
therefore $O(1)$ amortised.
\end{prop}

Equivalently, each element is pushed once, moved once and popped once: at most
three operations over its lifetime. Both stacks are preallocated at
initialisation so the ingest path performs no heap allocation, and the
accumulators are reset when the window size or sample clock is reconfigured,
since the aggregate is only meaningful relative to a fixed window definition.

% ─────────────────────────────────────────────────────────────────────────
\section{SEC --- Secrets and distribution}

\subsection{Key-based redaction severs the data dependency}
\label{lbl:HS-SEC-01}

Let $K$ be the set of secret key \emph{names} and $V$ the corresponding values.
Value-based scrubbing computes $\mathrm{redact}(\mathrm{text}, V)$ --- it must
hold the plaintext to match against it, so the logging subsystem becomes a
place the secret exists, and the matched values persist in string buffers.
Key-based redaction computes $\mathrm{redact}(\mathrm{env}, K)$, splitting each
entry at the first \texttt{=} and replacing the value when the key is in $K$.

\begin{prop}
$\mathrm{redact}(\cdot, K)$ has no data dependency on $V$: its output is
constant in the secret values. The logging path therefore never holds
plaintext, and the key set is not itself sensitive.
\end{prop}

This is an information-flow argument, and it is why the rule is ``scrub with
keys, not values'' rather than ``scrub carefully''.

\paragraph{Merge order.} The never-override-explicit rule states that an
explicit caller value wins over an injected default. Under Python's mapping
union the \emph{right} operand wins, so the correct expression is
\begin{equation}
\mathtt{secret\_env} \mid \mathtt{caller\_env},
\end{equation}
and the form \texttt{caller\_env | secret\_env} implements the opposite of the
stated rule. The narrative pattern document writes the latter; the register
carries the corrected form.

\paragraph{The AST guard is a guard, not a proof.} Scanning for forbidden path
literals and forbidden writer names is a syntactic over-approximation of the
property ``this module never writes a secret to disk''. It is sound for the
literals it enumerates and unsound in general, since a path assembled by
concatenation or read from configuration escapes it. It must therefore inspect
both \texttt{FunctionDef} and \texttt{AsyncFunctionDef} --- distinct AST nodes,
and a guard that checks only the former has a hole exactly the size of the
codebase's async surface --- and it must be understood as raising the cost of
the mistake rather than excluding it.

\paragraph{Fail closed on decryption.} A missing key or corrupt envelope
terminates the process. Proceeding with a partially populated configuration
runs the system with secrets silently absent, which is the
\ref{lbl:HS-BOUND-01} failure: an unsafe path that looks like a successful one.
Where a consumer genuinely requires a filesystem path, the render target must be
an \emph{asserted} tmpfs (verified in \texttt{/proc/mounts}, not assumed from
the path name), created \texttt{0600} at creation time rather than chmod'ed
afterwards, and unlinked on exit.

\subsection{Renaming bindings is \texorpdfstring{$\alpha$}{alpha}-conversion;
  renaming attributes is not}
\label{lbl:HS-SEC-02}

Let $B$ be the set of names the package binds and $K \subseteq B$ the keep set.
The applied map is
\begin{equation}
\alpha(x) = \begin{cases}
\mathrm{hash}(x, \mathrm{salt}) & x \in B \setminus K \\
x & \text{otherwise,}
\end{cases}
\end{equation}
required to be injective on $B \setminus K$ and disjoint from $K$, or two
identifiers merge and the program changes meaning.

\begin{prop}
Renaming binding occurrences is a capture-avoiding $\alpha$-conversion, and
is therefore semantics-preserving. Renaming an attribute access on a foreign
object is not an $\alpha$-conversion, and is unsound.
\end{prop}

\begin{proof}
$\alpha$-conversion renames a binder together with the occurrences it binds;
the term's meaning is invariant because the binder's identity is not observable
outside its scope. An attribute access \texttt{o.environ} is a field selection
on a value the package did not bind: the name belongs to \texttt{o}'s type, not
to this package's scope, so no binder is being renamed and the selection simply
fails to resolve after the rewrite.
\end{proof}

This is the precise reason for the scope rules: foreign attributes, string
literals, numeric literals and the \emph{external} names of keyword arguments
are all outside $B$. The one case needing special handling is \texttt{import x},
which binds \texttt{x} in the module namespace while naming an external module;
it must become \texttt{import x as \_alias} so the binding follows the rename
while the module reference does not.

\paragraph{The keep set is the API.} $K$ is an explicit, mechanically enforced
statement of the published surface: everything in $B \setminus K$ is private by
construction, and the build proves it. This is the most valuable artifact the
pipeline produces.

\paragraph{Determinism.} The build is a function of $(\mathrm{source},
\mathrm{salt}, \mathrm{toolchain})$. Pinning all three makes the artifact a
deterministic function of reviewed inputs, so a reviewer can rebuild with the
release salt and compare hashes; ``build twice, compare'' belongs in the build's
own test suite. The salt makes the mapping a build parameter rather than a
constant, so a release and an experiment are distinguishable without editing
source.

\paragraph{The driver is a program because two of its steps read artifacts.}
Asserting that no kept symbol was renamed (before compiling) and that no source
file shipped (after packaging) both require inspecting a produced artifact,
which a command sequence cannot do without becoming a program. A driver that
continues past a failed rewrite ships plain source, so each stage must abort the
build.

\paragraph{Scope limit, stated plainly.} Compilation to native extensions raises
the cost of reading the code; it is not encryption. The artifact still contains
a complete description of the semantics, so nothing that must remain secret ---
keys, credentials --- may be embedded, regardless of the compilation.

% ─────────────────────────────────────────────────────────────────────────
\section{KNOW --- Retrieval}

\subsection{Two incomparable orders, composed}
\label{lbl:HS-KNOW-01}

Vector retrieval ranks by similarity in an embedding of \emph{descriptions};
graph retrieval ranks by adjacency in a relation over \emph{structure}. Neither
order refines the other: two functions may be semantically similar and
structurally unrelated (two independent implementations of the same idea), and
two functions may be structurally adjacent and described in disjoint vocabulary
(a caller and its callee in different subsystems). Because the orders are
incomparable, neither retrieval alone is complete for the other's queries ---
which is the formal content of ``vectors locate, graphs explain''.

The pipeline composes them rather than choosing: the vector phase supplies
recall over description space by producing anchors from natural language, and
the graph phase supplies closure over structure by expanding each anchor along
\texttt{CALLS}, \texttt{CALLED\_BY}, \texttt{INHERITS} and \texttt{IMPORTS}.

\paragraph{The hop cap is a growth bound.} A neighbourhood of radius $h$ under
branching factor $b$ has size $O(b^h)$. Since the result must fit a context
window, $h$ must be capped --- $h \le 2$ in practice. An uncapped expansion
does not degrade gradually; it exhausts the window on one anchor's call tree.

\paragraph{Dual-write atomicity.} If the two stores are written in separate,
independently failing stages, the graph can hold a symbol whose chunk the vector
store lacks, and expansion returns a dangling reference. Ingest must therefore
write both within one transaction or one sync stage. De-duplication of identical
source spans happens before prompt assembly, since an anchor and its own
expansion routinely name the same span and paying for it twice is a direct loss
of context budget.

% ─────────────────────────────────────────────────────────────────────────
\section{Provenance}

The 34 rules above were folded on 2026-08-29 from the long-form pattern
documents held in \texttt{\textasciitilde/Docs/}, one file per pattern. Those
documents remain the narrative source, and each is named per rule in the
\texttt{pattern} field of the TOML SSOT. Two defects in those documents were found during the fold and are
corrected here rather than propagated: the secret-injection merge order
(\ref{lbl:HS-SEC-01}) and the short-name derivation rule
(\ref{lbl:HS-BOUND-06}). Both corrections are recorded in the register's
\texttt{invariant} fields, and both narrative documents should be amended to
match.

\end{document}
